% ── Geographic Taxonomy ───────────────────────────────────────────────────────
% Edit this file to customise the "Geographic Taxonomy" section intro.

\noindent
Geographic units are identified by \textit{extended GISCO IDs}, a harmonised coding
scheme derived from Eurostat's NUTS.
Extended GISCO IDs can take three forms, corresponding to different geographic levels:
\begin{itemize}
\item \textbf{Country codes} are two-letter codes aligned with European Union
conventions, which follow
\href{https://www.iso.org/obp/ui/#search/code/}{ISO 3166-1 alpha-2}
except for Greece (\texttt{EL}) and Kosovo (\texttt{XK}).
This corresponds to the \texttt{NUTS 0} level.
\item \textbf{Regional GISCO IDs} take the form
\texttt{CCL} (\texttt{NUTS 1}),
\texttt{CCLL} (\texttt{NUTS 2}),
or \texttt{CCLLL} (\texttt{NUTS 3}),
with \texttt{CC} a country code and each \texttt{L} alphanumeric.
Prefixes identify parent units, e.g.\
$\mathtt{PT10}$ is a subunit of $\mathtt{PT1}$, which is a subunit of
$\mathtt{PT}$ (Portugal).
\item \textbf{Local GISCO IDs} take the form \texttt{CC\_L\ldots L}.
For standard LAU municipalities the local part matches the Eurostat LAU code
(e.g.\ \texttt{AT\_10101} for Vienna's first district).
Extended local IDs of the form \texttt{CC\_CCYYxx} or \texttt{CC\_CCZZxx}
identify special aggregate units described in the subsections below.
Municipality-level units constitute the \texttt{LAU 2} level; identifiable
groups of municipalities are coded at \texttt{LAU 1}; other local units
(constituencies, postal aggregates, etc.) have a special type.
\end{itemize}

The subsections below describe how BLUE\_DB extends the standard GISCO taxonomy
to accommodate geographic units not present in the official NUTS/LAU
classification.
